\subsection{Profundidad de Anidamiento}

\subsubsection{Regla: máximo tres niveles de anidamiento}

\textbf{Regla:} No se anidan más de tres niveles de estructuras de control por método. Cuenta como nivel toda estructura que abre un bloque: \texttt{if}, \texttt{else}, \texttt{for}, \texttt{foreach}, \texttt{while}, \texttt{do}, \texttt{switch}, \texttt{try}, \texttt{catch}, \texttt{finally}, \texttt{using} y \texttt{lock}. El cuerpo del método es el nivel 0; la primera estructura dentro de él es el nivel 1. Un \texttt{else} o un \texttt{catch} ocupa el mismo nivel que su \texttt{if} o su \texttt{try}. Cuando se supera el límite, se combinan condiciones o se extrae la parte interna a un método privado (McConnell, 2004).

\textbf{Código correcto:}
\begin{lstlisting}[style=csharpstyle]
public void ProcessTurn(Player player)
{
    if (player is null || !player.IsActive)
    {
        return;
    }

    ExecuteTurn(player);
}

private void ExecuteTurn(Player player)
{
    try
    {
        _board.MovePlayer(player);
    }
    catch (InvalidMoveException ex)
    {
        _logger.Warn($"Move rejected. Player={player.Name}.", ex);
    }
}
\end{lstlisting}

\textbf{Código incorrecto:}
\begin{lstlisting}[style=csharpstyle]
public void ProcessTurn(Player player)
{
    if (player is not null)
    {
        if (player.IsActive)
        {
            try
            {
                if (_board.HasSpaceFor(player))
                {
                    _board.MovePlayer(player);
                }
            }
            catch (InvalidMoveException ex)
            {
                _logger.Warn($"Move rejected. Player={player.Name}.", ex);
            }
        }
    }
}
\end{lstlisting}
